\documentclass[10pt,a4paper]{amsart}
\usepackage[margin=19mm]{geometry}
\usepackage{amsmath,amssymb,mathtools}
\usepackage[scr=rsfs,cal=euler]{mathalfa}
\usepackage{xfrac}
\usepackage[unicode,hidelinks,bookmarks=false]{hyperref}
\newcommand{\ie}{i.e.\ }
\newcommand{\cf}{cf.\ }
\newcommand{\resp}{respectively\ }
\newcommand{\Subsection}{\subsection}
\providecommand{\nobreakdash}{\nobreak\hbox{-}}
\setlength{\emergencystretch}{2em}
\multlinegap0pt
\usepackage{amssymb}
\usepackage{dsfont}
\usepackage{mathtools}
\usepackage{xcolor}
\usepackage{float}
\newtheorem{assumption}{Assumption}
\newtheorem{lemma}{Lemma}
\newcommand\EE{\mathbb E}
\newcommand\RR{\mathbb R}
\newcommand\cA{\mathcal A}
\newcommand\cF{\mathcal F}
\newcommand\cG{\mathcal G}
\newcommand\cL{\mathcal L}
\newcommand\cS{\mathcal S}
\title{Markov Decision Processes of the Third Kind: Learning Distributions by Policy Gradient Descent}
\author{Nicole B\"auerle, Athanasios Vasileiadis$^*$}
\date{Source: arXiv:2602.06567v1 (CC BY 4.0)}
\begin{document}
\maketitle

\noindent\emph{Source and licence.} Markov Decision Processes of the Third Kind: Learning Distributions by Policy Gradient Descent. Version arXiv:2602.06567v1 (CC BY 4.0), distributed by arXiv under the Creative Commons Attribution 4.0 International licence, http://creativecommons.org/licenses/by/4.0/, as recorded in the arXiv licence metadata for this exact version and shown on the abstract page https://arxiv.org/abs/2602.06567v1. Full licence notice: \texttt{licenses/arxiv-2602.06567-CC-BY-4.0.txt}.\par\smallskip

\noindent\textit{Editorial scope.} Counted unit: the article's lemma lem:bias (source0.tex lines 730--761). The article's complete original proof of it is delivered with it (source0.tex lines 1776--1912, one \texttt{proof} environment of 137 lines, which the article places away from the statement). No mathematics is written, repaired or re-derived: the statement and the proof are the article's own lines, in the article's own notation, and no retained line is substituted or re-flowed.  Retained uncounted as the source-local context the proof needs: lines 518--526; lines 633--639; lines 1555--1560.  Nothing was omitted from any retained span, so the package carries no omission record.  No retained line contains a citation, so the wrapper carries no bibliography and no bibliography entry.  No retained line contains a figure environment, an image-inclusion command or drawing code, so no image is delivered and the package declares no asset dependency.  Editorial formatting only, in addition: an AMS \texttt{amsart} preamble that restates the article's own declarations and macros used by the retained lines, namely the environments \texttt{assumption}, \texttt{lemma} on separate counters, together with the standard packages those lines need.  The retained environments print in this document as Assumption 1, Lemma 1 in order of appearance, whereas the article prints the same items with the numbering of its own full document.  The complete native source of the article as distributed in the arXiv e-print for this exact version is bundled unchanged as \texttt{sources/194107/source0.tex}.
\begin{assumption} \label{ass:rFLipgrad}
We assume
\begin{itemize}
\item[(i)]$\partial_s r, \partial_a r$ exist and are  Lipschitz continuous, i.e. for all $s,\tilde s\in \cS,a, \tilde a \in \cA $ and for $i=s,a$
$$|\partial_i r(s,a)-\partial_ir(\tilde s,\tilde a)| \le L_{\nabla}^{r,i} \big( |s-\tilde s| + |a-\tilde a|  \big)$$
\item[(ii)] $\partial_s  F, \partial_a F$ exist and are Lipschitz continuous, i.e. for all $s,\tilde s\in \cS,a, \tilde a \in \cA,\varepsilon,\tilde\varepsilon\in\RR $ for $ i=s,a$
$$|\partial_iF(s,a,\varepsilon)-\partial_i F(\tilde s,\tilde a,\tilde\varepsilon)| \le L_{\nabla}^{F,i} \big( |s-\tilde s| + |a-\tilde a|\big). $$
\end{itemize}
\end{assumption}

 \begin{equation}\label{eq:approxgradient}
   g(\theta)
:= \nabla_\theta \hat\cL_{L}(\theta)= 2\sum_{\ell=1}^{L} \beta_\ell
\mathrm{Re}\Big(
 \overline{(\varphi^*(u_\ell)-\hat\varphi_\theta(u_\ell))} \cdot \widehat{\nabla_\theta \varphi_{\theta}}(u_\ell)
\Big).  
 \end{equation}

\begin{lemma}\label{lem:bound1}
    Under  Assumption \ref{ass:rFLipgrad} we obtain constants $B_\nabla^{s}, B_\nabla^a, B_\nabla^R$ such that on $\Theta\times G$ for all $t$:
    \begin{equation}
        \|\nabla_\theta s_t^\theta(\omega)\| \le B_\nabla^{s},\;  \|\nabla_\theta a_t^\theta(\omega)\| \le B_\nabla^a,\;  \|\nabla_\theta R_t^\theta(\omega)\| \le B_\nabla^R
    \end{equation}
\end{lemma} 

\begin{lemma}\label{lem:bias}
    The bias of the estimator for the gradient of the target function in \eqref{eq:approxgradient} at iteration time $k$ is given by
\begin{align}
\label{eq:bias-decomposition}
\EE\!\big[g(\theta_k)\mid \cF_k^G\big]
=
\nabla_\theta \cL_{L}(\theta_k)
+\mathrm{Bias}_M(\theta_k),
\end{align}
where
\begin{align}
\label{eq:bias-term}
\mathrm{Bias}_M(\theta_k)
=
\frac{2}{M}\sum_{\ell=1}^{L}\beta_\ell\,
\text{Re}\Big(
iu_\ell\,\EE[\nabla_\theta R_T^{\theta_k}|\cF_k^G]
-\overline{\varphi_{\theta_k}(u_\ell)}\,\nabla_\theta \varphi_{\theta_k}(u_\ell)
\Big),
\end{align}
and it satisfies the bound ($B_\nabla^R$ is the constant from Lemma \ref{lem:bound1} )
\begin{equation}
\label{eq:bias-bound}
\big\|\mathrm{Bias}_M(\theta_k)\big\|_2
\;\le\;
\frac{\hat{S}}{M} , \quad \hat{S}:=4\,B_\nabla^R\Big(\sum_{\ell=1}^{L}|\beta_\ell u_\ell|\Big).
\end{equation}
Moreover, we have
\begin{equation}
\label{eq:square-bound} \EE\!\big[ \|g(\theta_k)\|_2^2\mid \cF_k^G\big]\le \hat{S}^2.%16 (B_\nabla^R)^2 \Big(\sum_{\ell =1}^L |\beta_\ell u_\ell|\Big)^2
\end{equation}
\end{lemma}

\begin{proof}
of Lemma \ref{lem:bias}: We fix an iteration index $k\in\mathbb N$ and work throughout \emph{conditionally} on the
$\sigma$--algebra $\cF_k^G$.
Under this conditioning, the parameter $\theta_k$ is deterministic, and all bounds derived
on the good event $G$ hold pathwise.
In particular, by Lemma \ref{lem:bound1}
\[
\|\nabla_\theta R_T^{\theta_k}\|\;\le\; B_\nabla^R
\qquad\text{almost surely on } G.
\]

%For each frequency node $\omega_\ell$, 
Let $(R_m,\nabla_\theta R_m)_{m=1}^M$ be i.i.d.\ copies of
$(R_T^{\theta_k},\nabla_\theta R_T^{\theta_k})$ under the current parameter $\theta_k$. Thus, these Monte Carlo samples are conditionally i.i.d.\ given $\cF_k^G$. All expectations below are therefore taken with respect to these samples only, while $\theta_k$ is treated as fixed.

\medskip
\noindent\textbf{Step 1: Bias.}

Fix a frequency node $u=u_\ell$ and, to simplify notation, suppress the index $\ell$
throughout this step.
Recall that for this node the empirical characteristic function is given by
\[
\widehat\varphi(u)
:=
\frac1M\sum_{j=1}^M e^{iu R_j}.
\]

For each sample index $j\in\{1,\dots,M\}$, define the pathwise gradient contribution
\[
U_j(u)
:=
iu\,e^{iu R_j}\,\nabla_\theta R_j
\;\in\;\mathbb{C}^7.
\]
With this notation, the contribution of the frequency node $u_\ell$ to the full estimator
$g(\theta_k)$ can be written as
\[
g_\ell
=
\frac{2\beta_\ell}{M}\sum_{j=1}^M
\text{Re}\Big(
\overline{\widehat\varphi(u)-\varphi^\star(u)}\;
U_j(u)
\Big).
\]

We emphasize that the same Monte Carlo batch
$(R_m)_{m=1}^M$ is used both to construct the empirical characteristic function
$\widehat\varphi(u)$ and to evaluate the gradient terms $U_j(u)$.
This coupling is the source of the finite-sample bias analyzed below. Taking conditional expectation and using linearity of $\text{Re}(\cdot)$,
\[
\EE[g_\ell\mid \cF_k^G]
=
\frac{2\beta_\ell}{M}\sum_{j=1}^M
\text{Re}\Big(\EE[\overline{\widehat\varphi(u)}\,U_j(u)\mid \cF_k^G]
-\overline{\varphi^\star(u)}\,\EE[U_j(u)\mid \cF_k^G]\Big).
\]
Because the samples are i.i.d., it suffices to compute $\EE[\overline{\widehat\varphi(u)}\,U_1(u)\mid \cF_k^G]$:
\[
\EE[\overline{\widehat\varphi(u)}\,U_1(u)\mid \cF_k^G]
=
\frac1M\EE[\overline{e^{iu R_1}}\,U_1(u)\mid \cF_k^G]
+\frac{M-1}{M}\EE[\overline{e^{iu R_2}}\mid \cF_k^G]\EE[U_1(u)\mid \cF_k^G].
\]
Now $\overline{e^{iu R_1}}\,U_1(u)= iu \nabla_\theta R_1$,
and we denote $\EE[{e^{iu R_2}}\mid \cF_k^G]=\EE_G[{e^{iu R_T^{\theta_k}}}] =:{\varphi_{\theta_k}(u)}$.
Thus $\EE[U_1(u)\mid \cF_k^G]=\nabla_\theta\varphi_{\theta_k}(u)$
and we obtain
\[
\EE[\overline{\widehat\varphi(u)}\,U_1(u)\mid \cF_k^G]
=
\frac1M\,iu\,\EE[\nabla_\theta R_1\mid \cF_k^G]
+\frac{M-1}{M}\,\overline{\varphi_{\theta_k}(u)}\,\nabla_\theta\varphi_{\theta_k}(u).
\]
Plugging back yields
\[
\EE[g_\ell\mid \cF_k^G]
=
2\beta_\ell\,\text{Re}\Big(\overline{\varphi_{\theta_k}(u)-\varphi^\star(u)}\,\nabla_\theta\varphi_{\theta_k}(u)\Big)
+\frac{2\beta_\ell}{M}\text{Re}\Big(iu\,\EE[\nabla_\theta R_1\mid \cF_k^G]-\overline{\varphi_{\theta_k}(u)}\,\nabla_\theta\varphi_{\theta_k}(u)\Big).
\]
Summing over $\ell$ gives \eqref{eq:bias-decomposition}--\eqref{eq:bias-term}.

To bound the bias, use $\|\nabla_\theta R_1\|_2\le \|\nabla_\theta R_1\|\le B_\nabla^R$ and
$\|\nabla_\theta\varphi_{\theta_k}(u)\|_2
\le \|\nabla_\theta\varphi_{\theta_k}(u)\|
\le \EE[|u|\,\|\nabla_\theta R_1\|\mid \cF_k^G ]\le |u| B_\nabla^R$,
and $|\varphi_{\theta_k}(u)|\le 1$, obtaining \eqref{eq:bias-bound}.

\medskip
\noindent\textbf{Step 2: Second moment.}
For each $(\ell,j)$ define
\[
Y_{\ell,j}
:=\mathrm{Re}\Big(\overline{\hat{\varphi}(u_\ell)-\varphi^*(u_\ell)}\,iu_\ell e^{iu_\ell R_j}\nabla_\theta R_j\Big)\in\RR^7,
\qquad
g(\theta_k)=\frac{2}{M}\sum_{\ell=1}^{L}\beta_\ell\sum_{j=1}^M Y_{\ell,j}.
\]
We have
$|\hat{\varphi}(u_\ell)-\varphi^*(u_\ell)|\le  2$
and $|e^{i\cdot}|=1$. Hence, for all $(\ell,j)$ we obtain on  $G$
\begin{equation}\label{eq:Y-bound}
\|Y_{\ell,j}\|
\le 2|u_\ell|\,\|\nabla_\theta R_j\|
\le 2|u_\ell|\,B_\nabla^R.
\end{equation}

Next, by the triangle inequality,
\[
\Big\|\sum_{j=1}^M Y_{\ell,j}\Big\|
\le \sum_{j=1}^M \|Y_{\ell,j}\|
\le M \max_{1\le j\le M}\|Y_{\ell,j}\|.
\]
%Squaring and using \eqref{eq:Y-bound} yields the deterministic estimate, valid on $\cG_\delta^{(T)}$,
%\begin{equation}\label{eq:sumj-bound}
%\Big\|\sum_{j=1}^M Y_{\ell,j}\Big\|^2
%\le M^2 \big(2|u_\ell|K_{\mathrm{grad}}^{(\delta)}\big)^2
%= 4M^2u_\ell^2\big(K_{\mathrm{grad}}^{(\delta)}\big)^2.
%\end{equation}
%Plugging \eqref{eq:sumj-bound} into the definition of $g(\theta_k)$ and using again the triangle inequality,
Thus, we obatin
\[
\|g(\theta_k)\|
\le \frac{2}{M}\sum_{\ell=1}^{L}|\beta_\ell|
\Big\|\sum_{j=1}^M Y_{\ell,j}\Big\|
\le \frac{2}{M}\sum_{\ell=1}^{L}|\beta_\ell|\,
\Big(2M|u_\ell|B_\nabla^R\Big)
=4B_\nabla^R\sum_{\ell=1}^{L}|\beta_\ell u_\ell|.
\]
Therefore, on $G$,
\begin{equation}\label{eq:g-det-bound}
\|g(\theta_k)\|_2^2
\le 16\big(B_\nabla^R\big)^2
\Big(\sum_{\ell=1}^{L}|\beta_\ell u_\ell|\Big)^2.
\end{equation}
Taking conditional expectations preserves the inequality, so we obtain the statement.
\end{proof}

\end{document}
